\chapter{The Noncomputable Core}
\label{chap:noncomputable-core}

% Closed question: What selects one admissible continuation as the history
% actually carried forward?

\begin{chapteressay}

This is the chapter the book has been building toward.

The first five chapters developed the algorithmic vocabulary:
enumeration, divide and conquer, interpolation, union-find,
bisection. Each chapter named a closed question and showed how the
admissible record forced a specific family of algorithms. The
algorithms in each chapter were computable: they could be
implemented in Lean, run on a CPU, profiled, optimized, and
certified by a typeclass instance. Every algorithm produced its
output in finite time.

This chapter is different. The closed question is: what selects one
admissible continuation as the history actually carried forward? The
answer the book argues for is: a selection rule based on minimum
Kolmogorov complexity. The argument that follows is the load-
bearing argument of the book. The selection rule requires exact
Kolmogorov complexity. Exact Kolmogorov complexity is provably
uncomputable. Therefore the selection rule is uncomputable. The
universe, if it is running the selection rule, is doing something
no Turing machine can simulate.

The argument is not metaphorical. It is technical and the
techniques are decades old. Turing proved the Halting Problem
uncomputable in 1936. Chaitin proved the noncomputability of
Kolmogorov complexity in the 1960s. The Busy Beaver function was
defined by Rado in 1962 and shown to grow faster than any computable
function. These are established results in computer science. The
book applies them to the question of physical selection.

The chapter's first concrete result is the Halting Problem. There
is no algorithm that, given a program, decides whether the
program halts on all inputs. Turing's diagonal argument is the
proof. The proof is short and exactly the kind of argument that
formalist mathematics produces. The Halting Problem is the
foundational uncomputability result; everything else in this
chapter reduces to it.

The chapter's second concrete result is the Busy Beaver function.
$BB(n)$ is the maximum number of steps any halting $n$-state
Turing machine performs. It is the chapter's hard wall: $BB$ grows
faster than any computable function, so no algorithm can keep up
with it. $BB(5) = 47{,}176{,}870$ was proved in 2024 by exhaustive
analysis. $BB(6)$ is unknown. The function's uncomputability is
what prevents the computation of exact Kolmogorov complexity. To
compute $K(w)$ for a string of length $\ell$, the algorithm would
need to time each candidate program by $BB$(its length); since $BB$
is uncomputable, the algorithm cannot do this.

The chapter's third concrete result is Chaitin's $\Omega$, the
halting probability. $\Omega$ is a specific real number between 0
and 1, defined by summing $2^{-|p|}$ over all halting programs.
The binary digits of $\Omega$ are algorithmically random and
encode the Halting Problem in compact form. The first $n$ digits
of $\Omega$ would suffice to decide halting for all programs of
length up to $n$. The digits exist; they are not computable.

These three results --- the Halting Problem, Busy Beaver,
Chaitin's $\Omega$ --- are the chapter's technical core. They
establish that the selection rule is uncomputable. The remaining
two sections (heartbeat proxies in Lean,
\texttt{ArmWaveProcess} as overhead meter) examine what
algorithms can do given the uncomputability: they can use
proxies that approximate the selection rule. The proxies are
useful but not exact.

The chapter's second concrete example is the sealed archive. A
government archive holds records of varying restriction levels.
Some boxes are sealed; their contents are restricted. The
archivist works with a catalog that identifies each box and its
restriction level. She routes requests through the catalog. She
does not read the contents of sealed boxes; her clearance does
not extend that far.

The catalog is a proxy for the archive's content. It supports
routing decisions; it does not contain the underlying records.
Two boxes with very different contents may have identical
catalog entries (if both are sealed). The catalog admits the
existence of sealed boxes; it does not reveal what is inside.

This is the chapter's algorithmic structure in archival form.
The selection rule (the actual decision of who can access what)
depends on the catalog (the proxy). The catalog is computable;
the contents of sealed boxes are not. The archive operates by
using the proxy where the underlying content is inaccessible.

The chapter's closed question is:

\begin{center}
\emph{What selects one admissible continuation as the history
actually carried forward?}
\end{center}

The answer the chapter develops, the load-bearing answer of the
book, is: the selection requires exact Kolmogorov complexity,
which is uncomputable. The universe somehow has access to
exact $K$; no Turing machine does. The selection rule is the
distinguishing feature between the universe's algorithmic
behavior and any simulable algorithm.

The chapter's argument is not pessimistic. The proxies are
useful. The compiler's heartbeat counter is a working
approximation of computational cost. The \texttt{ArmWaveProcess}
cascade is a working approximation of formal overhead. These
proxies allow algorithms to make admissibility decisions even
when exact $K$ is unavailable. The work proceeds, despite the
fundamental limitation.

The chapter's contribution is to draw the line precisely. The
proxies are computable; the selection rule is not. The proxies
are good for many purposes; they are not exact. The chapter
develops the boundary with concrete examples: the Halting
Problem (Turing 1936), the Busy Beaver function (Rado 1962),
Chaitin's $\Omega$ (Chaitin 1975). Each is a piece of established
computer science. The book applies them to the question of
physical selection.

The reader who reaches this chapter is invited to consider its
implications carefully. If the selection rule is uncomputable,
then no simulation of the universe on a Turing machine can be
faithful. The simulation might be very close --- using
sophisticated proxies, using extensive computation, using
clever approximations --- but it cannot be exact. The universe's
behavior at the level of selection involves a noncomputable
operation. The simulation hypothesis, which posits that the
universe is running on a Turing-like machine, is at odds with
this conclusion.

The book does not claim mysticism. The selection rule is a
mathematical object. The argument is technical. The
conclusion is that the universe is doing something different in
kind from algorithmic computation. The difference is precise:
the universe has access to exact Kolmogorov complexity, which is
provably noncomputable. No Turing machine does. The mathematics
is the same; the implementation is different.

\end{chapteressay}

\section{The Halting Problem}
\label{se:the-halting-problem}

A computer program either halts (produces an output and stops) or runs
forever (never produces a final output). The question is: given an
arbitrary program $p$ and an arbitrary input $x$, can an algorithm
decide which?

In 1936, Alan Turing proved that no such algorithm exists. The proof
is by diagonalization. Suppose an algorithm $H(p, x)$ exists that
returns true if $p$ halts on input $x$, and false otherwise. Construct
a new program $D$ as follows: given input $p$, compute $H(p, p)$. If
$H(p, p) = $ true (meaning $p$ halts when given itself as input), $D$
enters an infinite loop. If $H(p, p) = $ false, $D$ halts.

Now consider $H(D, D)$. Two cases.

Case one: $H(D, D) = $ true, meaning $D$ halts on input $D$. By the
construction of $D$, if $H(D, D) = $ true, then $D$ enters an infinite
loop on input $D$. Contradiction: $D$ both halts and does not halt.

Case two: $H(D, D) = $ false, meaning $D$ does not halt on input $D$.
By the construction of $D$, if $H(D, D) = $ false, then $D$ halts on
input $D$. Contradiction.

Either way, $H$ produces a contradiction. The assumption that $H$
exists is false. There is no algorithm that decides halting in
general.

The argument is short and elementary. It is also profound. It says
that the question ``does this program halt?'' is not algorithmically
decidable. The question is well-formed --- given a specific program
and input, the answer is either yes or no --- but no algorithm can
answer it for all inputs.

The Halting Problem is the canonical example of an undecidable
problem. Many other problems are undecidable too, often by reduction
to halting: the equivalence of two arbitrary programs, the
non-emptiness of an arbitrary recursively enumerable set, the
validity of arbitrary first-order logic sentences, the convergence
of an arbitrary numerical iteration. Each of these problems contains
the Halting Problem as a special case; if any of them were
decidable, halting would be too, which contradicts Turing.

For the chapter's argument, the Halting Problem is the foundation of
the noncomputability of Kolmogorov complexity. The Kolmogorov
complexity of a string $w$ is the length of the shortest program that
produces $w$ and halts. To compute $K(w)$, the algorithm must
identify the shortest halting program for $w$. To do that, it must
test programs in order of increasing length and check which ones
halt on the empty input producing $w$. The check is the Halting
Problem. If any algorithm could compute $K$, it could also decide
halting --- which is impossible.

The reduction is direct. Given a program $p$, define $w_p$ to be
some encoding of $p$'s behavior. If $K(w_p)$ were computable, the
algorithm computing it would have to determine which programs halt
producing $w_p$, which means deciding halting for those programs.
The Halting Problem reduces to Kolmogorov complexity computation.
Therefore Kolmogorov complexity is uncomputable.

The Halting Problem is the bottleneck. Every noncomputability
result in this chapter traces back to it. The Busy Beaver function
is uncomputable because it grows faster than any computable
function. Chaitin's $\Omega$ is uncomputable because its bits
encode halting probabilities. Kolmogorov complexity is uncomputable
because it requires finding shortest halting programs. The
selection rule that picks the minimum-complexity admissible
continuation is uncomputable because it requires Kolmogorov
complexity.

The Lean compiler encounters the Halting Problem at every
elaboration step. When a definition is recursive, the compiler must
verify termination: that the recursion eventually halts. The
verification uses a well-founded relation --- a function from the
recursion's state to natural numbers, decreasing at each call,
bounded below by zero. The well-founded relation is the user's
proof that the recursion terminates. Without the proof, the
compiler refuses to admit the definition.

This is the compiler's accommodation of the Halting Problem.
Termination is undecidable in general, so the compiler does not
attempt to decide it. Instead, the compiler asks the user to
provide a termination proof, and the compiler checks the proof. The
proof is a finite object; the check is a finite procedure. The
undecidability is sidestepped by requiring the user to supply
enough information to make the decision finite.

This is also the chapter's first concrete example of the
algorithm-versus-universe distinction. The user (or the
mathematician, or the physicist) provides a termination proof
based on their understanding of the recursion's structure. The
compiler verifies the proof. The user's understanding is something
the compiler cannot derive automatically. The compiler can only
check what the user provides. If the user cannot provide a
termination proof (because the function does not terminate), the
compiler refuses to admit it.

The Halting Problem thus sets up the chapter's argument. The
selection rule that picks the minimum-Kolmogorov-complexity
continuation requires solving the Halting Problem. The Halting
Problem is provably unsolvable by any algorithm. Therefore the
selection rule is not computable. The chapter develops the
consequences in the remaining sections.

\section{Busy Beaver}
\label{se:busy-beaver}

The Busy Beaver function $BB(n)$ is defined as follows. Consider all
Turing machines with $n$ states (plus a halt state), an alphabet of two
symbols (0 and 1), and a one-dimensional tape initially filled with
zeros. Run each machine starting from a designated start state. Some
machines halt; some do not. Of the machines that halt, identify the one
that performs the most steps before halting. $BB(n)$ is the number of
steps performed by that maximum machine.

The definition is constructive: for any specific $n$, the answer is a
specific natural number. For $n = 1$, $BB(1) = 1$. For $n = 2$,
$BB(2) = 6$. For $n = 3$, $BB(3) = 21$. For $n = 4$, $BB(4) = 107$. For
$n = 5$, the value was finally pinned down in 2024: $BB(5) =
47{,}176{,}870$.

For $n = 6$, $BB(6)$ is unknown. Lower bounds have been pushed steadily
upward over the decades. As of recent work, $BB(6)$ is known to exceed
$10^{36{,}534}$, and the actual value may be vastly larger. The problem
of finding $BB(6)$ exactly is at the limit of current proof-theoretic
techniques.

The Busy Beaver function grows faster than any computable function. The
proof: suppose $f$ is computable and $f(n) \ge BB(n)$ for all
sufficiently large $n$. Construct a Turing machine $M_f$ that takes
input $n$ and computes $f(n)$. The size of $M_f$ is some fixed
constant, say $k$ states. For inputs $n > k$, the machine $M_f$ takes
at most $f(n)$ steps. But $M_f$ run on input $n$ is a $k$-state
machine (the size of $M_f$ does not depend on the input), and the
Busy Beaver function $BB(k)$ is the maximum steps for any $k$-state
machine that halts. So $f(n) \le BB(k)$ for all $n$, contradicting
the assumption that $f(n) \ge BB(n)$ for large $n$. The function $f$
cannot dominate $BB$.

The argument shows that $BB$ grows so fast that no computable
function can keep up with it. The Busy Beaver function is
uncomputable. There is no algorithm that, given $n$, produces $BB(n)$
for all $n$. Specific values can be computed (for small $n$) by
exhaustive analysis of all $n$-state machines, but the analysis itself
requires solving halting for each machine, and the halting problem is
undecidable. For small $n$, the analysis is feasible; for larger $n$,
it requires increasingly sophisticated techniques, and beyond some
threshold, the techniques themselves run out.

The Busy Beaver function is the chapter's hard wall.

To see why it is the wall, consider the chapter's argument about the
selection rule. The selection rule requires picking, from among all
admissible continuations of a record, the one with minimum Kolmogorov
complexity. To do this, the algorithm must enumerate candidate
continuations, compute their complexities, and select the minimum.
The complexity of a candidate $w$ is $K(w)$, the length of the
shortest halting program that produces $w$.

To compute $K(w)$, the algorithm tests programs in order of
increasing length: programs of length 1, then 2, then 3, and so on.
For each program, it runs the program and checks whether the output
is $w$. If the program produces $w$ and halts, the algorithm has
found a candidate. The minimum-length program that produces $w$ and
halts is the witness of $K(w)$.

The catch: for each program tested, the algorithm must decide
whether the program halts. If the program halts, the algorithm
observes the output and proceeds. If the program does not halt, the
algorithm waits forever. There is no a priori bound on the running
time of a halting program; some halting programs take an enormous
number of steps before producing their output. The algorithm cannot
simply time out, because timing out might prematurely declare a
slow-halting program as non-halting, missing a candidate.

This is where the Busy Beaver function enters. Suppose the algorithm
wants to compute $K(w)$ where $w$ has length $\ell$. The programs to
test have length up to roughly $\ell$ (longer programs are
guaranteed not to be the minimum). Each program of length $L$ has
internal state space bounded by $2^L$ (each bit of the program is
either 0 or 1). When the program runs on the empty input, its
computation visits at most $2^L \cdot \ell$ tape states before
repeating (since the relevant tape positions are bounded by $\ell$,
the length of the output).

If the program is going to halt, it will halt within $BB(L)$ steps.
If the program runs for more than $BB(L)$ steps without halting, it
will run forever. So the algorithm could, in principle, test each
length-$L$ program by running it for exactly $BB(L)$ steps and
checking whether it halted by then. If it halted, observe the
output. If not, declare it non-halting.

The catch within the catch: this requires knowing $BB(L)$. And
$BB(L)$ is uncomputable. The algorithm needs the value of an
uncomputable function. The algorithm cannot have it. Therefore the
algorithm cannot complete the Kolmogorov complexity computation. The
selection rule based on Kolmogorov complexity is not algorithmic.

The Busy Beaver function is the precise reason. If $BB$ were
computable, the algorithm could time each candidate program for
$BB(L)$ steps, identify the halting ones, find the shortest, and
return $K(w)$. The selection rule would be computable. Since $BB$
is not computable, the algorithm cannot do this. The selection rule
is not computable.

This is the deepest result of the book. The argument is concrete.
The selection rule requires $K(w)$. The computation of $K(w)$
requires $BB(L)$. The Busy Beaver $BB$ is uncomputable. Therefore
the selection rule is uncomputable.

Many readers will object that approximate $K$ might suffice for
many purposes. The chapter's response is: approximate $K$ is
computable (Lempel-Ziv, gzip, and other compression algorithms
provide upper bounds on $K$, and these bounds are computable). But
the chapter's argument is about the exact $K$, not approximate.
Exact $K$ is what the selection rule needs, because approximate $K$
can sometimes be wrong by more than the difference between two
candidates. Two candidates with similar approximate $K$ may have
very different exact $K$; the selection between them depends on
which has the truly minimum.

For most practical purposes, approximate $K$ works because most
candidates are clearly different in complexity. The hard cases are
near-ties: two candidates whose Kolmogorov complexities differ by
one or two bits. In near-ties, approximate $K$ cannot resolve. The
exact $K$ would, but exact $K$ is not computable. Near-ties are
inadmissible to the algorithm; they are admissible only to whatever
entity has access to exact $K$.

The Busy Beaver function therefore is the algorithmic boundary
between admissible computation and inadmissible computation. The
boundary is real. It is not a matter of insufficient computational
power. No amount of additional CPU, memory, or time can move the
boundary. The boundary is set by the structure of computation
itself.

What about proxy costs? Several proxies are computable. The
heartbeat counter in Lean is one. The CPU time of a running
program is another. The memory footprint is a third. These
proxies correlate with $K$ for many practical cases. They diverge
from $K$ in cases of high asymmetry: a program with low Kolmogorov
complexity but high runtime cost (e.g., one that computes $BB(100)$
exactly --- short to write, eternal to run), or a program with
high Kolmogorov complexity but low runtime cost (e.g., one that
contains a large hard-coded constant and just prints it).

The proxies are useful for many practical purposes. They are not
the selection rule. The selection rule requires exact $K$. The
proxies provide approximate $K$. The two are not the same.

A reader skeptical of the book's claim should focus on this gap.
The claim is not that the universe runs faster than your laptop.
The claim is that the selection rule requires a quantity (exact
$K$) that no Turing machine can compute. The universe, if it is
running the selection rule, is doing something Turing machines
cannot. The Busy Beaver function is the exact reason.

For larger Turing machines --- machines with 7 or more states ---
the Busy Beaver value is so large that, even if it could somehow
be computed, the resulting time to enumerate halting programs of
the corresponding lengths would exceed the age of the universe by
absurd factors. $BB(7)$ is widely believed to be G\"odel-like in
its uncomputability: not just unreachable by current methods, but
provably independent of the axioms of standard set theory (ZFC).
The book does not require this stronger claim; it suffices that
$BB$ is not computable.

The chapter's closing thought on this section: $BB(5) =
47{,}176{,}870$ was proved in 2024. $BB(6)$ is unknown. The
progress has been steady, but it is finite progress against an
infinite mountain. The wall does not move. The wall is what the
chapter has identified.

\section{Chaitin's Omega}
\label{se:chaitins-omega}

Chaitin's constant $\Omega$ is the probability that a randomly generated
program halts. The construction is precise. Fix a universal Turing
machine $U$ that takes a prefix-free input encoding. For each input $p$
of length $|p|$, the probability of generating $p$ by flipping a fair
coin is $2^{-|p|}$. The sum over all halting inputs is:
\[
\Omega = \sum_{p : U(p) \text{ halts}} 2^{-|p|}.
\]
The prefix-free property ensures that $\Omega$ is at most 1; it is also
positive (some programs halt) and less than 1 (some programs do not).
$\Omega$ is a specific real number in $(0, 1)$.

The remarkable fact is that the binary digits of $\Omega$ are
algorithmically random. The first $n$ binary digits of $\Omega$, taken
as a string, are incompressible: no program shorter than $n$ bits can
produce them. The Kolmogorov complexity of the first $n$ digits is
approximately $n$ itself. The digits are random in Chaitin's strong
sense: they are random not just statistically, but algorithmically ---
no algorithm can compute them in shorter form than reading them.

The digits encode the halting problem in a compact form. The first $n$
digits of $\Omega$ are enough to decide, in principle, whether each
program of length at most $n$ halts. Knowing the first $n$ digits, one
could enumerate all programs of length up to $n$, run them in
parallel, and as halting programs are observed, accumulate their
contribution to $\Omega$. When the accumulation matches the first $n$
digits to within $2^{-n}$ precision, the algorithm has identified all
the halting programs of length up to $n$. The remaining programs do
not halt.

This means the first $n$ digits of $\Omega$ are an oracle for the
Halting Problem restricted to programs of length up to $n$. The
digits are an oracle for a uncomputable problem. They are not
themselves computable; computing them would solve the Halting Problem.

$\Omega$ is the algorithmic version of the Holy Grail. It is a
specific number, mathematically well-defined, that solves an
uncomputable problem. Its existence is provable; its value is
uncomputable. No algorithm can produce more than finitely many of
its digits.

For the chapter's argument, $\Omega$ is the most precise expression
of the boundary. The selection rule requires identifying admissible
continuations and computing their Kolmogorov complexities. The
identification involves halting tests; the halting tests are
oracular in $\Omega$. If the algorithm had access to $\Omega$'s
digits, it could perform the halting tests and the Kolmogorov
complexity computations. Without $\Omega$, it cannot.

The universe, if it implements the selection rule, has implicit
access to $\Omega$ (or to some equivalent oracle). The universe
does what the algorithm cannot do, because the universe has access
to the uncomputable information that the algorithm lacks.

In Lean, $\Omega$ appears as the \texttt{ChaitinsNumberSequence}
typeclass introduced in Chapter 3. The typeclass carries the digits
as a sequence of Booleans, with the constraint that the $n$-th
Boolean encodes the halting status of the $n$-th program. The type
is well-formed; no instance is constructively buildable.

The compiler can prove theorems about $\Omega$ --- that it is
algorithmically random, that it is uncomputable, that its digits
encode the Halting Problem. The compiler cannot evaluate any
specific digit. The asymmetry between proof and computation is
sharp: the theorem about $\Omega$ is checkable in finite time; the
digit of $\Omega$ is not computable in finite time.

This connects $\Omega$ to the book's larger argument about type-
checking versus execution. The Lean type-checker can admit
statements about $\Omega$ as well-formed propositions. The Lean
runtime cannot evaluate $\Omega$. The two are different
operations on the same object: checking the type is finite;
evaluating the value is not.

The chapter's algorithmic claim is that $\Omega$ is the canonical
example of an object that the type system admits but the runtime
cannot construct. The selection rule, which depends on quantities
like $\Omega$, falls into the same category. The selection rule is
type-correct (it has a well-formed mathematical specification); it
is not runtime-computable.

This is the gap the book identifies. Type-correct does not imply
runtime-computable. The mathematical universe admits well-formed
objects whose values cannot be computed. The physical universe, if
it is implementing the selection rule, is constructing values of
such objects --- doing what no Turing machine can do.

$\Omega$ also appears in connections to algorithmic randomness and
information theory. Specifically, $\Omega$ is the canonical example
of an incompressible infinite sequence: the first $n$ bits cannot be
generated by any program shorter than $n$ bits. The sequence has
maximum entropy in the algorithmic sense. The chapter's next
sections develop these connections: proxy costs in Lean, formal
overhead as a measure of admissibility, and the distinction
between approximate complexity (computable) and exact complexity
(not computable).

\section{Heartbeat Proxy in Lean}
\label{se:heartbeat-proxy-in-lean}

Lean's elaborator has a built-in mechanism for bounding its own
running time. The setting \texttt{set\_option maxHeartbeats 400000}
caps the elaboration of a single declaration at 400,000 heartbeats. A
heartbeat is a unit of internal work --- roughly, one elementary
unification step or one term reduction. When the elaborator reaches
the cap, it stops and reports a timeout. The declaration fails.

The heartbeat is the elaborator's proxy for the Kolmogorov
complexity of the declaration. Definitions with high heartbeat counts
tend to be complex: they involve deep typeclass resolution, many
implicit arguments, or large term reductions. Definitions with low
heartbeat counts tend to be simple: a numeric literal, a
straightforward function definition, a direct application of a
basic operator.

The correlation is imperfect. The heartbeat count depends on
implementation details: the order in which constraints are gathered,
the specific simp rules invoked, the caching state of the
elaborator. Two definitions with the same Kolmogorov complexity may
have different heartbeat counts. Two definitions with different
Kolmogorov complexities may have the same heartbeat count.

The proxy is useful despite its imperfection. The compiler uses it
to prevent infinite elaboration: if the elaborator is using too many
heartbeats, the declaration is rejected, regardless of whether it
would eventually succeed. The heartbeat limit is the compiler's
defense against pathological inputs.

The \texttt{EKG} type in \texttt{LeanCalibration.lean} exposes the
heartbeat counter to user code:

\begin{leanexcerpt}{EKG}
private structure Reading where
  heartbeats : Nat

structure EKG where
  private reference : Reading
  embigger? : Prop $\to$ Prop $\to$ Prop
\end{leanexcerpt}

The \texttt{Reading} structure is private: user code cannot
directly read the heartbeat count. The \texttt{EKG} structure
exposes the \texttt{embigger?} relation, which takes two
propositions and returns a proposition asserting that the first
dominates the second under the calibration's bookkeeping. The
\texttt{reference} field is itself private: user code cannot
directly read the baseline either. The comparison through
\texttt{embigger?} is the public surface; both the count and the
baseline are hidden implementation.

This design is intentional. The cost of computations is exposed as
\texttt{Prop}-valued statements, not as runtime numbers. ``Program
A costs less than program B'' is a proposition. The proposition
can be proved or disproved; it can be used in type signatures and
theorems; it can constrain other propositions. The
\texttt{Reading.heartbeats} field is the underlying mechanism, but
the public interface treats cost as a propositional comparison, not
as a measurable quantity.

This is the chapter's algorithmic version of calibration as hidden
reference. The heartbeat counter is the reference scale --- the
calibration's hidden mechanism. The public interface exposes only
the propositional comparisons --- the calibrated readings. The
private/public split is the algorithm's commitment that costs are
admissible only through the calibrated interface.

The \texttt{LOGICAL} typeclass carries the \texttt{EKG}:

\begin{leanexcerpt}{LOGICAL}
class LOGICAL
    (Value : Type)
    (Carrier : CarrierProcess Value)
    [DISTINGUISHABLE Value Carrier] ... [UNIVERSAL Value Carrier]
  where
  feelings : HeartbeatProcess Value Carrier
  ekg      : Calibration.EKG
  logical? : YarnTheory $\to$ YarnTheory $\to$ Prop := fun a b =>
    YarnTheory.le a b
\end{leanexcerpt}

A \texttt{LOGICAL} instance is a program that carries its own
calibration. The calibration is the \texttt{EKG}; the
\texttt{feelings} field names the underlying
\texttt{HeartbeatProcess}; the \texttt{logical?} method is the
binary decision rule over \texttt{YarnTheory} that the program's
costs are measured by. The cascade through \texttt{UNIVERSAL}
carries every prior class as a typeclass requirement.

The chapter's argument is that the heartbeat is the proxy because
exact $K$ is not available. The proxy is computable; exact $K$ is
not. The selection rule, in the universe's implementation, uses
exact $K$. The compiler's selection (which definitions to admit,
which to reject) uses the heartbeat proxy. The proxy is good
enough for the compiler's purposes; it is not good enough for the
selection rule the chapter is identifying as noncomputable.

The proxy makes a useful upper bound. If two programs have very
different heartbeat counts, the program with the lower count
probably has lower Kolmogorov complexity. The probably is the
proxy's limitation: in pathological cases, the proxy's prediction
is wrong. But for typical programs, the proxy correlates well with
$K$, and the compiler's selection (based on the proxy) approximates
the selection rule (based on $K$).

The book's claim is that the universe is doing what the compiler
cannot: using exact $K$ instead of the proxy. The chapter
demonstrates the compiler's proxy and notes its limitations. The
limitations are not contingent; they are structural. No proxy can
do what exact $K$ does, because exact $K$ is not computable.

\section{ArmWaveProcess as Overhead Meter}
\label{se:armwaveprocess-as-overhead-meter}

\texttt{Episode7.lean} introduces a cluster of typeclasses that track
formal overhead accumulated during elaboration. The cluster is named
satirically: \texttt{ArmWaveProcess} for the process,
\texttt{BULLSHIT} for the certificate, \texttt{Diatribe},
\texttt{CrusadeProcess}, \texttt{PROPAGANDA}, \texttt{Cult},
\texttt{InitiationProcess}, and \texttt{ACOLYTE} for the cascade.

The satirical names are deliberate. The typeclasses measure something
real (formal overhead, the accumulated burden of typeclass resolution
and term construction during elaboration), but the measurement is
proxy, not exact $K$. The satire is a reminder that the proxy is
not the truth --- it is the gestures the compiler makes while doing
the work.

\texttt{ArmWaveProcess}:

\begin{leanexcerpt}{ArmWaveProcess}
structure ArmWaveProcess
    (Value : Type)
    (Carrier : CarrierProcess Value)
    [DISTINGUISHABLE Value Carrier] ... [FINITE_ELEPHANT Value Carrier]
  where
  galerkin_process : GalerkinProcess Value Carrier
  guess            : Spline
  reticulate?      : Spline $\to$ Spline := ...
\end{leanexcerpt}

An \texttt{ArmWaveProcess} carries the full project cascade
through \texttt{FINITE\_ELEPHANT}, a current \texttt{Spline} guess,
and a \texttt{reticulate?} function that advances the guess one
arm-wave step. The cascade is the proof obligation that the
arm-wave sits on every prior class admissibly. The structure
formalizes the observation that elaboration is, in part, a
sequence of formal gestures (the spline reticulations) that
contribute to the proof but do not necessarily contribute to the
content.

\texttt{BULLSHIT}:

\begin{leanexcerpt}{BULLSHIT}
class BULLSHIT
    (Value : Type)
    (Carrier : CarrierProcess Value)
    [DISTINGUISHABLE Value Carrier] ... [FINITE_ELEPHANT Value Carrier]
  where
  arm_wave_process : ArmWaveProcess Value Carrier
  interpolate?     : Spline $\to$ Spline $\to$ Prop := ...
\end{leanexcerpt}

A process is \texttt{BULLSHIT} when it carries an
\texttt{ArmWaveProcess} together with a binary decision rule
\texttt{interpolate?} on \texttt{Spline}s. The decision rule
checks whether two splines admissibly interpolate one to the
other; the cascade obligation runs the entire length of the prior
class stack. The certificate is the conservation principle the
chapter names: bullshit is strictly conserved under reticulation.

The cascade continues with \texttt{Diatribe} (a longer-than-needed
argument), \texttt{CrusadeProcess} (a process that propagates a
particular \texttt{Diatribe} pattern), \texttt{PROPAGANDA} (the
certificate that the propagation has occurred), \texttt{Cult} (a
group of definitions that all share a \texttt{PROPAGANDA}), \texttt{
InitiationProcess} (the process by which a new definition enters
the \texttt{Cult}), and \texttt{ACOLYTE} (the certificate that a
definition has been initiated).

Each step in the cascade adds formal overhead. A \texttt{Diatribe}
has more overhead than an \texttt{ArmWaveProcess}; a
\texttt{Cult} has more overhead than a \texttt{Diatribe}; an
\texttt{ACOLYTE} has the most overhead in the cascade. The
typeclasses measure progressive accumulation: each level adds
specific structure that the previous levels did not have.

The chapter's algorithmic claim is that this accumulation is the
compiler's running estimate of Kolmogorov complexity. The
estimate is by stages: each stage adds a fixed amount of
formal structure, and the total is the sum. The estimate is
computable (the typeclass instances can be checked in finite
time), but it is approximate, not exact.

The satirical naming serves a purpose. The cascade reminds the
reader that the formal overhead is bureaucratic. It is the
accumulated formal apparatus around what may be a simple
content. The chapter's argument is that this bureaucratic
proxy is the best the compiler can do; the true complexity is
not the proxy but the underlying content, and that true
complexity is what no algorithm can compute exactly.

The book's larger argument benefits from this satire because it
makes the proxy's role explicit. The compiler has a proxy. The
proxy is the cascade of typeclasses. The proxy has a cost
(formal overhead). The cost is computable. The compiler's
selection (which definitions to admit) is based on the proxy.
The universe's selection (which physical history to actualize)
is based on something the proxy cannot capture.

The book's contention is that the proxy can be very useful for
many practical purposes but is fundamentally different from
exact $K$. The cascade in \texttt{Episode7.lean} is the
compiler's working proxy. The proxy can be measured, compared,
optimized, and reasoned about. The exact $K$ cannot. The two
are not the same.

This concludes Chapter 6, the load-bearing chapter of the book.
The argument has been: the selection rule requires exact $K$;
exact $K$ requires solving the Halting Problem (via Busy
Beaver); the Halting Problem is provably undecidable; therefore
the selection rule is uncomputable. The compiler has proxies
(heartbeats, the \texttt{ArmWaveProcess} cascade) that
approximate $K$. The proxies are useful but they are not exact
$K$. The universe, implementing the selection rule, has access
to exact $K$ in some form. No Turing machine can.

\begin{bridge}

The chapter's question was what selects one admissible continuation
as the history actually carried forward.

The answer: the selection rule picks the admissible continuation
with minimum exact Kolmogorov complexity. The Halting Problem is
the foundational uncomputability result. The Busy Beaver function
is the hard wall: it grows faster than any computable function, and
its uncomputability prevents the timing of halting tests required
to compute exact $K$. Chaitin's $\Omega$ encodes the halting
problem in a single uncomputable real number; its digits would
provide an oracle if accessible, but they are not computable. The
heartbeat proxy in Lean and the \texttt{ArmWaveProcess} typeclass
cascade are the compiler's approximate substitutes for $K$.
Useful approximations; not exact $K$.

What remains is transport. The selection rule is uncomputable, but
the universe somehow implements it. Information moves between
admissible records, transported by some structure the algorithm
can describe but not simulate. Chapter 7 introduces inversion ---
the algorithmic form of moving information backward through a
computation, from outputs to inputs.

\end{bridge}

\begin{coda}{The Sealed Archive}

The archive occupies the basement of a government building. Steel
shelving runs in rows under fluorescent lights. Climate control
maintains the temperature at sixty-eight degrees and the humidity at
forty percent. Each shelf holds rows of acid-free archival boxes,
labeled with accession numbers and date ranges. The archive holds
records related to certain government activities; many of the
boxes are sealed, with their contents restricted.

The archivist sits at a desk near the entrance. Her job is to
respond to access requests from researchers, government agencies,
and occasionally journalists. Each request specifies an accession
number; the archivist looks up the box in the catalog, determines
whether the requester has clearance for that box's restriction
level, and either retrieves the box or denies the request. The
archivist herself does not read the contents of sealed boxes; her
clearance does not extend to the interior of most of the archive.

The catalog is on a computer terminal beside her desk. Each entry
lists the accession number, the acquisition date, the donor or
source, a one-line description, and a restriction level. The
restriction level determines who can request the box. Some boxes
are public: the description tells what is inside. Some boxes are
restricted but described: the description names the topic without
revealing the content. Some boxes are sealed: the description
reads ``[CONTENT RESTRICTED]'' and the topic itself is classified.

The catalog has thirty-four thousand entries. Of these, about
five thousand are sealed. The archivist can route requests through
the catalog quickly; she can determine, for any accession number,
the restriction level and whether the requester has clearance. She
cannot determine, from the catalog alone, what is in any sealed
box. The catalog is a proxy: it carries enough information to
route, not enough information to assess.

A researcher arrives this morning with a request. He is a
historian working on a book about a thirty-year-old policy
decision. He has obtained clearance from the relevant agency to
view three specific accession numbers. The archivist verifies the
clearance, locates the boxes, and brings them to the reading room.
The researcher spends the day with the boxes.

At the end of the day, the researcher returns the boxes. The
archivist re-shelves them. She notes the researcher's visit in
the access log. She has performed her routing function correctly:
the right boxes were retrieved, the right researcher was granted
access, the activity was logged. She did not learn what was in
the boxes; she did not need to.

Now consider a different request. A young academic comes in with a
list of fifteen accession numbers. The list looks plausible, but
the academic does not have clearance for any of them. The
archivist checks each number against the clearance database. Each
returns the same result: insufficient clearance. The archivist
denies the requests. She notes the denials in the log. The
academic leaves.

Did the denials reveal anything about the content of the requested
boxes? Strictly speaking, no. The denials say only that the
academic does not have clearance. They do not say what the boxes
contain. The catalog descriptions, which the academic could read
publicly, listed all fifteen as ``[CONTENT RESTRICTED]''. The
denials did not change that.

But there is a subtle point. The denials confirmed that the
boxes exist. The academic now knows that there are fifteen
specific accession numbers with restricted content. Each number
is a position in the catalog. Each position has been verified to
exist. The verification is a small piece of information: the
existence of the box. The verification is not the content.

The catalog is the algorithmic proxy. It carries accession
numbers (the names of records), restriction levels (the routing
information), and descriptions (the public content, where
non-restricted). It does not carry the sealed content. The proxy
is useful for routing; it does not carry the underlying truth.

The chapter's argument lives in this distinction. The proxy
allows admissibility decisions: which boxes can be retrieved,
which cannot, which descriptions can be shared. The proxy does
not provide the actual content. The actual content is in the
boxes; the boxes are sealed; the seals are the algorithmic
boundary.

Some boxes have additional restrictions. A subset of the
sealed boxes carry an annotation: ``[CONTENT
RESTRICTED --- UNVERIFIED]''. These are boxes whose contents
have never been verified by anyone with clearance. They were
sealed at their creation and have remained sealed. The catalog
entry lists them, the access log records that no researcher
has ever viewed them, the descriptions are blank. The boxes
exist; their content is unknown to every living person.

These are the chapter's deepest analog. The catalog admits the
existence of the boxes. The boxes have specific positions on
the shelves. The boxes have weights, sizes, manufacturing
dates. The shelving is physically real. But the content is
unverified. No one has read it. No one knows what it says.

To read these boxes would require breaking the seal --- an
action that itself is restricted, requiring clearance levels
that no living person currently has. The seals are the
algorithmic Halting Problem. The content is the
\texttt{ChaitinsNumberSequence}: a well-defined sequence of
bits that no algorithm can compute. The catalog is the
heartbeat proxy: a working approximation that supports the
routing operations but does not give access to the content.

The archivist closes the archive at five o'clock. She locks
the entry door, sets the security alarm, and walks to the
metro station. The catalog is updated with the day's
activity: three boxes retrieved, three returned, fifteen
requests denied, no seals broken. The state of the archive is
the same at five as it was at nine, plus the audit trail of
the day's transactions.

The archive's records include the
\texttt{ArmWaveProcess} cascade in legal form. Every request,
denied or granted, is logged. The log is the bureaucratic
overhead: the formal record of what was attempted, what was
permitted, what was denied. The overhead is real; it consumes
filing cabinet space, archivist's time, computer memory in the
audit database. The overhead is the proxy for the content of
the archive: it does not contain the content, but it tracks
the access patterns around the content.

A researcher attempting to game the archive --- to figure out
what is inside the sealed boxes from the access pattern alone
--- would face a problem analogous to the chapter's
noncomputability argument. The access patterns are proxy
information. They can hint at the content; they cannot
reveal it. Two boxes with very different content might have
identical access patterns. Two boxes with identical content
might have very different access patterns. The proxy is
useful for routing; it is not exact $K$ of the content.

This is the chapter's final image. The archive is the
universe. The boxes are the admissible continuations. Some
are accessible; some are sealed; some have content
unverifiable by any algorithm. The catalog is the proxy;
the heartbeat counter is the proxy; the
\texttt{ArmWaveProcess} cascade is the proxy. The proxies
support the work of routing, of admitting, of denying. The
proxies are not the content of the boxes.

The content of the universe's admissible continuations
exists. Some entity --- the universe itself --- has access
to the exact $K$ of each continuation. That entity
implements the selection rule. The catalog of the
universe (the proxy) supports the admissibility decisions
the algorithm can make. The seals of the universe (the
exact $K$) support the decisions the algorithm cannot.

The archivist takes the metro home. Tomorrow she will return.
The archive will be exactly as she left it, plus whatever
additions arrive overnight. The boxes will remain on the
shelves. The seals will remain sealed. The catalog will
remain a proxy. The work will continue.

The catalog is computable. The boxes are not.

\end{coda}
